% 中文摘要

物联网与人工智能技术的持续发展推动了虚拟现实、增强现实、自动驾驶等新兴应用的落地，同时也对网络边缘的计算能力提出了更高要求。边缘计算通过将计算任务卸载至网络边缘节点执行，有效降低了响应时延。然而，随着应用数据的激增和计算任务规模的扩大，单一节点难以独立支撑大规模计算任务，分布式边缘计算因此受到广泛关注。分布式边缘计算通过聚合多个边缘节点的计算资源，将大规模任务分解为多个子任务并行处理，从而提升计算效率。但由于边缘节点的硬件配置、负载状态和网络环境存在异构性，部分节点可能因资源竞争或突发故障成为迟滞节点，其计算速度远落后于其他节点，导致整体计算时延受限于最慢节点。为缓解迟滞效应，编码边缘计算作为一种融合纠删码理论与分布式计算的新兴范式被提出，通过在任务分发阶段引入冗余，使中心节点仅需收集部分完成较快节点的计算结果即可重构原始任务结果。然而，现有编码边缘计算研究大多聚焦于计算时延优化，忽视了下行传输阶段由于信道衰落可能导致的通信时延。事实上，编码计算在任务分发阶段引入的冗余使得各边缘节点的计算结果之间存在强相关性，这为利用分布式信源编码（Distributed Source Coding，DSC）技术优化下行传输提供了可能。针对上述问题，本文从DSC的视角出发，分别针对完整恢复计算结果和线性恢复计算结果两种场景，研究了编码计算结果传输方案与DSC的联合优化方法。本文的主要研究内容与贡献如下：

（1）针对需要完整恢复子任务计算结果的编码计算场景，提出了一种面向完整恢复的DSC辅助编码计算结果传输方案。该方案利用编码计算结果间的线性相关性，结合Slepian-Wolf（SW）定理，给出了节点间无通信条件下实现计算结果完整恢复的理论压缩区域。由于SW定理仅适用于理论研究，实际系统中难以实现，设计了基于伴随式穿刺的编码方案，使边缘节点能够根据信道状态自适应调整压缩率，实现传输负载在节点间的灵活重分配。基于所提方案，建立了下行传输时延最小化优化模型。针对该非凸优化问题，设计了基于凹凸过程（Convex-Concave Procedure，CCCP）的联合优化算法，联合优化各边缘节点的信源压缩率与传输功率。仿真结果表明，所提方案在信道条件较差场景下降低了传输时延，相比独立编码方案具有明显的性能提升，且对迟滞节点数量变化具有良好的鲁棒性。

（2）针对仅需恢复计算结果线性函数值的编码计算场景（如分布式机器学习中的梯度聚合），提出了一种面向线性恢复的DSC辅助编码计算结果传输方案。该方案结合嵌套线性码，给出了仅需要恢复计算结果线性函数值场景下的可达信源压缩区域；并证明了在特定条件下，面向线性恢复的可达速率区域严格大于经典SW区域，从而实现了比完整恢复更高的压缩效率。基于所提方案，给出了可达速率区域的简化形式，即SW边界与线性恢复边界的并集，降低了优化问题的求解复杂度；并进一步建立了相应的时延最小化优化模型，设计了基于CCCP的优化算法。仿真结果表明，所提方案的下行传输性能均优于独立编码方案和SW编码方案，在迟滞节点较多的场景下优势更为明显。

综上所述，本文通过挖掘编码边缘计算中任务间的强相关性，将DSC与编码计算相结合，为解决边缘计算系统中的通信瓶颈问题提供了新的视角与有效方法。